% Generated draft: requires mathematical and source review.
\documentclass[11pt]{article}
\usepackage[margin=25mm]{geometry}
\usepackage{fontspec}
\setmainfont{Latin Modern Roman}
\usepackage{amsmath,amssymb,mathtools,unicode-math}
\setmathfont{Latin Modern Math}
\usepackage{enumitem,needspace,xurl,hyperref}
\setlength{\parindent}{0pt}
\setlength{\parskip}{5pt}
\setlength{\emergencystretch}{3em}
\newcommand{\N}{\mathbb{N}}
\newcommand{\Z}{\mathbb{Z}}
\newcommand{\Q}{\mathbb{Q}}
\newcommand{\R}{\mathbb{R}}
\newcommand{\C}{\mathbb{C}}
\newcommand{\F}{\mathbb{F}}
\newcommand{\A}{\mathbb{A}}
\newcommand{\PP}{\mathbb{P}}
\newcommand{\E}{\mathbb{E}}
\newcommand{\eps}{\varepsilon}
\newcommand{\dd}{\,\mathrm{d}}
\begin{document}
\section*{UnsolvedMath ID: 30001245}
\textbf{Source problem code:} catalogue entry
\subsection{Definitions and mathematical statement}
To make the supplied question precise under conventional interpretations, let \(k\) be an integer with \(k\ge 2\). A finite \(k\)-uniform hypergraph is a pair \(H=(V(H),E(H))\), where \(E(H)\subseteq\binom{V(H)}{k}\). Its maximum vertex degree is
\[
\Delta(H)=\max_{v\in V(H)}\bigl|\{e\in E(H):v\in e\}\bigr|,
\]
with \(\Delta(H)=0\) if \(V(H)=\varnothing\). Its weak chromatic number \(\chi(H)\) is the least positive integer \(q\) for which there is a map \(\varphi:V(H)\to\{1,\ldots,q\}\) such that no edge is monochromatic. For a fixed \(k\)-uniform hypergraph \(F\), say that \(H\) is \(F\)-free if there is no injective map \(\iota:V(F)\to V(H)\) satisfying \(\iota(e)\in E(H)\) for every \(e\in E(F)\). Thus this interpretation excludes ordinary subhypergraph copies, without requiring them to be induced.

The supplied question, on the range where its displayed expression is defined, asks whether
\[
\begin{gathered}
\text{for every integer }k\ge2\text{ and every fixed finite }k\text{-uniform hypergraph }F,\\
\text{there exists }c_F>0\text{ such that every finite }F\text{-free }k\text{-uniform hypergraph }H\\
\text{with }\Delta(H)\ge2\text{ satisfies}\qquad
\chi(H)\le c_F\left(\frac{\Delta(H)}{\log\Delta(H)}\right)^{1/(k-1)}.
\end{gathered}
\]
Here \(\log\) denotes the natural logarithm, and \(c_F\) may depend only on \(F\), whose uniformity determines \(k\).

\textit{Unresolved qualifications.} The source does not explicitly specify finiteness, the range of \(k\), the copy convention, or the coloring convention; the definitions above describe a conventional reading requiring review. The supplied wording quantifies over every maximum degree, but its expression is undefined at \(\Delta=0\) and \(\Delta=1\). The displayed formulation therefore records only its meaningful range \(\Delta\ge2\); a small-degree convention or an explicitly asymptotic formulation remains to be supplied. No structural restriction on \(F\) beyond uniformity is stated, and none has been inserted.
\subsection{Short English statement}
The supplied question asks whether forbidding any fixed \(k\)-uniform hypergraph \(F\) forces the weak chromatic number to be at most a constant depending only on \(F\) times \((\Delta/\log\Delta)^{1/(k-1)}\). The draft spells out conventional finite, non-induced-copy interpretations; the intended conventions and treatment of degrees zero and one require review. The source labels the question open.
\subsection{Sources}
\begin{itemize}
\item Supplied local UnsolvedMath registry record, permanent ID \(30001245\), title ``Coloring \(F\)-Free Uniform Hypergraphs,'' source problem code \texttt{catalogue entry}, in the Oberwolfach Reports Open Problems collection. The record supplies proposed year \(2009\) but names no proposer.
\item Supplied report citation: \emph{Combinatorics and Probability} (2009), Oberwolfach Reports, DOI \texttt{10.4171/owr/2009/23}, pp.~1225--1302. \url{https://doi.org/10.4171/owr/2009/23}. The background describes these reports as documenting mathematical workshops and their open-problem discussions; no problem-specific page or named contributor is supplied.
\item Supplied literature reference: Bohman--Frieze--Mubayi, \emph{Coloring simple hypergraphs}. \url{https://doi.org/10.1016/j.jctb.2013.09.003}.
\item Supplied literature reference, labeled ``Frieze--Mubayi restricted theorem.'' \url{https://arxiv.org/abs/0809.2979}.
\item Literature assessment dated \texttt{2026-08-21}: the supplied curation labels the problem open and states that the conjectured bound is proved for important restricted classes, including ``simple hypergraphs,'' while the stated general \(F\)-free version remains conjectural. This is supplied dated curation, not an independently verified assessment. The background explicitly requests independent verification of time-sensitive status and cited claims.
\item The registry labels the statement \texttt{corrected\_verified} and describes its preparation as meaning-preserving semantic, typographic, and LaTeX cleanup with no unsupported source-dependent content added. These are supplied provenance claims, not mathematical verification of this draft.
\end{itemize}
\end{document}
